%  05_design.tex  –  Chapter 4: Design


\chapter{Design}
\label{ch:design}

\section{System Architecture}
\label{sec:system-arch}

The Indoor Air Station is composed of three main subsystems that work together to collect,
store, and present environmental data:

\begin{itemize}
  \item The Raspberry Pi Monitoring Station
  \item The AWS Cloud Backend
  \item The Website Frontend
\end{itemize}

Data transfer among the subsystems takes place using a simple linear pipeline process. The
sensors connected to the Raspberry Pi make measurements of the environment on a regular
interval of 30 minutes. The Raspberry Pi processes the unprocessed data from the sensors,
calculates the concentration of ppm for each gas sensor, compensates for temperature and
humidity changes, determines the presence of warnings, and finally writes the whole
measurement into the CSV local log file. After processing, the data is sent in a JSON payload
via HTTP to the Flask server located in AWS cloud. Once there, the Flask server will check the
validity of the data and add it to the PostgreSQL database present on AWS RDS service. If any
user wants to access the data, they must first visit the website where the Flask server makes a
query in the database and returns the results.

The website operates separately from the Raspberry Pi since even when the server is not
working, data collection and local logging go on without fail.

From the website it is available to download the Android mobile application. Similar to the
website, the app communicates with the Flask server over HTTP, consuming the same API
endpoints. No new server side components are added. It shows the data and predictions similar
to the website. It also contains the home and form pages.

\vspace{12pt}
\begin{figure}[H]
  \centering
  \includegraphics[width=0.85\textwidth]{system_architecture}
  \caption{Figure 10 - System architecture diagram}
  \label{fig:system-arch}
\end{figure}
\vspace{12pt}

\section{Software Architecture}
\label{sec:software-arch}

The software side of the project is divided into two distinct components: the station software
running on the Raspberry Pi and the backend running on AWS.

The station program consists of a Python program running all the time on the Raspberry Pi.
Initially, the software will set up the SPI for the MCP3208 ADC, initialize the DHT11
temperature and humidity sensor, setup the GPIO pins used by the LED, the button, and the
sound interrupt, get the geolocation of the station's device, and ask the user to input the
neighborhood name where the station is located. The software starts another thread which will
run the data collection process. A hardware interrupt function works separately from the other
threads by counting the number of times sound is sensed using the KY-037 digital output module.

The backend code comprises only one Flask app that processes the entire logic behind the server
operations. This application serves static pages and receives sensor information provided by
the Raspberry Pi, processes the authentication of user accounts by registering, verifying
email addresses, and logging in, resetting passwords, processing the contact forms'
information and serving filtered sensor information to the front-end pages and exporting the
information in CSV format for logged-in users. Database connections are established using the
connection pool technique to process simultaneous requests without establishing any new
connections.

\vspace{12pt}
\begin{figure}[H]
  \centering
  \includegraphics[width=0.85\textwidth]{software_architecture}
  \caption{Figure 11a - Software architecture diagram}
  \label{fig:software-arch}
\end{figure}
\vspace{12pt}

\begin{figure}[H]
  \centering
  \includegraphics[width=0.75\textwidth]{software_usecase}
  \caption{Figure 11b: Software architecture use case}
  \label{fig:software-usecase}
\end{figure}
\vspace{12pt}

\section{Hardware Architecture}
\label{sec:hardware-arch}

The hardware of the station consists of two breadboards positioned side by side next to the
Raspberry Pi 4B. The first breadboard has the DHT11 temperature and humidity sensor, the
MCP3208 ADC, the KY-037 sound sensor, the two LEDs and the push button. The second breadboard
carries the five MQ gas sensors: MQ135, MQ3, MQ6, MQ7, and MQ8. Both breadboards share a
common power and ground rail, supplied by the Raspberry Pi's 3.3\,V and 5\,V pins.

All five MQ sensors and the KY-037 analog output are connected to the MCP3208 ADC through
individual voltage dividers, each formed by a 1\,k$\Omega$ and 2\,k$\Omega$ resistor pair,
which scales the 5\,V sensor output down to the 3.3\,V range accepted by the ADC. The MCP3208
communicates with the Raspberry Pi over the SPI bus. The DHT11 is connected directly to a GPIO
pin with a 5\,k$\Omega$ pull-up resistor. The KY-037 digital output is connected directly to a
separate GPIO pin for interrupt-based sound event counting. The two LEDs each have a
220\,$\Omega$ current-limiting resistor. The push button is connected to a GPIO pin with the
internal pull-down resistor enabled in software.

The entire assembly is currently open hardware, though the compact form factor of the
components means it can be placed in a small container-like object, like a box for example.

\vspace{12pt}
\begin{figure}[H]
  \centering
  \includegraphics[width=0.85\textwidth]{hardware_architecture}
  \caption{Figure 12 - Hardware architecture diagram}
  \label{fig:hardware-arch}
\end{figure}
\vspace{12pt}
